% Generated draft: requires mathematical and source review.
\documentclass[11pt]{article}
\usepackage[margin=25mm]{geometry}
\usepackage{fontspec}
\setmainfont{Latin Modern Roman}
\usepackage{amsmath,amssymb,mathtools,unicode-math}
\setmathfont{Latin Modern Math}
\usepackage{enumitem,needspace,xurl,hyperref}
\setlength{\parindent}{0pt}
\setlength{\parskip}{5pt}
\setlength{\emergencystretch}{3em}
\newcommand{\N}{\mathbb{N}}
\newcommand{\Z}{\mathbb{Z}}
\newcommand{\Q}{\mathbb{Q}}
\newcommand{\R}{\mathbb{R}}
\newcommand{\C}{\mathbb{C}}
\newcommand{\F}{\mathbb{F}}
\newcommand{\A}{\mathbb{A}}
\newcommand{\PP}{\mathbb{P}}
\newcommand{\E}{\mathbb{E}}
\newcommand{\eps}{\varepsilon}
\newcommand{\dd}{\,\mathrm{d}}
\begin{document}
\section*{UnsolvedMath ID: 6200100}
\textbf{Source problem code:} catalogue entry
\subsection{Definitions and mathematical statement}
Let $\mathbb H^3$ be real hyperbolic three-space and identify its ideal boundary with $S^2$. The only explicit geometric data in the supplied item are subsets
\[
 \Lambda\subseteq S^2=\partial_\infty\mathbb H^3.
\]
No group is to be specified or assumed in the requested analogue.

The source asks for ``a similar statement to this inflexibility result'' for such subsets, but the referenced inflexibility theorem is absent from the record. Its assumptions, objects, deformations or maps, quantitative estimates, and conclusion are all missing. Consequently the source does not determine a mathematical implication that can be written with complete quantifiers. In particular it does not say whether $\Lambda$ must be closed, whether it is a limit set, or which regularity, measure, or complementary-domain conditions are allowed.

The task retained here is the source's request for a group-free subset version of that omitted theorem. It cannot be replaced by an arbitrary quasiconformal-extension or conformal-rigidity assertion: the research summary itself identifies those only as likely interpretations. Recovering the exact antecedent theorem is necessary to state the intended problem; no guessed theorem is supplied.
\subsection{Short English statement}
Can the inflexibility result referenced in the original source be generalized to subsets of the ideal two-sphere without choosing a group? The quoted theorem is missing, so the intended analogue is not yet a fully specified mathematical statement.

The supplied summary labels the question open but only guesses that it concerns quasiconformal extension or conformal rigidity. That guess is not adopted as the target. The 2005 source lists the problem, and provenance marks current status for review.
\subsection{Sources}
\begin{itemize}
\item Local UnsolvedMath registry, permanent ID 6200100, source code catalogue entry, ``Boundaries of Groups and Kleinian Groups --- Problem 100,'' proposed by Danny Calegari, 2005. Dataset curation and metadata: Ulam AI and UnsolvedMath Contributors, CC BY 4.0.
\item Kapovich, \emph{Problems on Boundaries of Groups and Kleinian Groups} (2005), Problem 100, PDF page 25, \url{https://www.math.ucdavis.edu/~kapovich/EPR/problems.pdf}; normalized pdf-text extraction, accessed 2026-07-29.
\item The antecedent inflexibility theorem and its reference are not present in the supplied record. The research summary supplies no verified replacement statement or later primary citation. Source-specific provenance marks current status and statement-reuse terms NEEDS\_REVIEW for the public author-hosted PDF, separately from the dataset licence.
\end{itemize}
\end{document}
